\cotefile{10}
\scanpage{171}
% archive sheet number: 170
topologique de $U$) dans $k$ (et $=$ sur $k$). Ainsi
on a un \deleted{\unsure{application}} homomorphisme d'algèbres
commutatives de $A$ dans l'algèbre des fonctions
\deleted{de} $G$ à valeurs dans $k$, transformant $1$ en $1$, compatible
d'ailleurs avec l'application diagonale et
l'augmentation, enfin aussi avec $\unsure{\nu|v|\vee}$. Dire que
c'est injectif signifie précisément que
$G$ est total dans $\underline U$. Noter que $G$ est exactement
l'ensemble de \uline{tous} les caractères de $A$
(qui est un groupe grâce à l'application
diagonale dans $A$ !) avec la topologie ``faible''
dire donc que $A\to\text{fonctions sur }G$ est
biunivoque signifie aussi que la transformation
de Gelfand \deleted{de} $A$ est biunivoque (si p. ex. $A$
est à \uline{engendrement fini sur $k$}, \deleted{\unsure{celà signifie que}}
\inserted{$k$ algébriquement}
% slanted
\margin{interpréter ceci en termes \unsure{d'associabilité|d'associativité} tensorielle $ {\unsure{K|k}}^{\underline U}$}
\uline{clos}, ceci signifie simplement que $A$ n'a pas
d'éléments nilpotents non nuls (si $x,\deleted{\unsure{\in|x'}}\in V$, $x'\in V'$
sont tels que $x'$ n'est pas orth. au plus
petit sous-\unsure{truc} de $V$ contenant $x$,
alors la $=$ ci-dessus est vraie pour
\[
x\otimes x\in V\otimes V\quad\text{et}\quad x'\otimes x'\in V'\otimes V'\,;
\]
contre-exemple \inserted{en car. $p\ne0$} à partir de $\underline A=k[X]/X^p\deleted{\unsure{-I}}$
\deleted{\unsure{$I$ puissance $p$ de l'idéal maximal|$I$ puissance pme de l'idéal maximal} $I=(X)^p$}
\deleted{$I$ l'idéal \unsure{maximal},} avec l'hom. diagonal naturel $\unsure{\leadsto}$ ; y-a-t-il un
\margin{???} contre-exemple en car. $0$ ?)
